% This file is part of the modular source for Stratified Relational Workspace Theory.
% !TeX root = ../Stratified_Relational_Workspace_Theory.tex
\theorypart{XIV}{Scaffolded and Endogenous Formation}

\topic{47}{Formation Systems and Scaffold Specification}

The preceding parts begin with a transfer system that has already been chosen.
A formation theory must instead explain how the nodes, edges, transfer maps,
workspace channels, and operational boundaries are obtained. This requires a
strict separation between the structure supplied by the experimenter and the
structure selected from interaction history.

\begin{definition}[Formation system]\label{def:formation-system}
A \emph{workspace formation system} is a tuple
\[
\mathfrak F
=
\bigl(
\Sigma,
\mathcal D_{\mathrm{form}},
\mathcal L,
\Omega,
\widehat\theta,
\mathcal D_{\mathrm{cert}},
\mathcal I_{\mathrm{form}}
\bigr),
\]
where:
\begin{itemize}
\item \(\Sigma\) is the externally supplied scaffold;
\item \(\mathcal D_{\mathrm{form}}\) is the formation or training history;
\item \(\mathcal L\) is the learning or selection rule;
\item \(\Omega\) is its internal random seed or stochastic state;
\item \(\widehat\theta=\mathcal L(\Sigma,\mathcal D_{\mathrm{form}},\Omega)\)
 is the learned organization;
\item \(\mathcal D_{\mathrm{cert}}\) is held-out certification data; and
\item \(\mathcal I_{\mathrm{form}}\) is a declared family of interventions on
 scaffold and formation history.
\end{itemize}
The learned organization is decomposed as
\[
\widehat\theta
=
\bigl(
\widehat V,
\widehat E,
\widehat H,
\widehat P,
\widehat{\mathcal O},
\widehat B
\bigr),
\]
where \(\widehat B\) is a learned operational boundary.
\end{definition}

The scaffold is itself stratified:
\[
\boxed{
\Sigma
=
\bigl(
\Sigma_V,
\Sigma_E,
\Sigma_H,
\Sigma_P,
\Sigma_{\mathcal O},
\Sigma_B,
\Sigma_Y
\bigr).
}
\]
The final coordinate \(\Sigma_Y\) records labels, targets, privileged masks,
or other supervision that could directly reveal the intended workspace. This
coordinate is necessary because an architecture may appear to discover a
boundary even when the desired membership mask was encoded in its training
labels.

\subtopic{47.1 Scaffolded architecture versus scaffolded solution}

A scaffold can constrain the search space without specifying its solution.
For example, a graph neural architecture can fix the class of admissible local
message-passing rules while leaving the active edges and transfer maps to be
learned. Conversely, a system can use a flexible architecture while receiving
labels that directly encode the desired workspace. The former is
architecturally scaffolded but solution-open. The latter is solution-scaffolded.

For any component \(q\) of \(\widehat\theta\), define its scaffold-only Bayes
risk under loss \(\ell_q\) by
\[
\mathcal R_q^{\Sigma}
=
\inf_f
\mathbb E
\left[
\ell_q\bigl(q,f(\Sigma,\Omega)\bigr)
\right].
\]
Define the full formation risk by
\[
\mathcal R_q^{\Sigma,\mathcal H}
=
\inf_g
\mathbb E
\left[
\ell_q\bigl(q,g(\Sigma,\mathcal D_{\mathrm{form}},\Omega)\bigr)
\right].
\]
The history advantage is
\[
\boxed{
\Delta_q^{\mathrm{hist}}
=
\mathcal R_q^{\Sigma}
-
\mathcal R_q^{\Sigma,\mathcal H}.
}
\]
A positive history advantage means that formation history contains information
about \(q\) that is not already recoverable from the declared scaffold and
random seed.

\begin{proposition}[Scaffold-only impossibility]\label{prop:scaffold-only-impossibility}
If
\[
q=f(\Sigma,\Omega)
\qquad\text{almost surely},
\]
then for every formation-history intervention that leaves \((\Sigma,\Omega)\)
fixed,
\[
\mathcal L\!\left(q\mid
\operatorname{do}(\mathcal D_{\mathrm{form}}=h)
\right)
=
\mathcal L\!\left(q\mid
\operatorname{do}(\mathcal D_{\mathrm{form}}=h')
\right),
\]
and \(\Delta_q^{\mathrm{hist}}=0\).
\end{proposition}

\begin{proof}
The value of \(q\) is measurable with respect to \((\Sigma,\Omega)\), so an
intervention on formation history cannot change its conditional law when those
variables are held fixed. The scaffold-only predictor \(f\) attains the same
risk as any predictor that also receives formation history, hence the risk gap
vanishes.
\end{proof}

This proposition is elementary but methodologically decisive. A structure
that is deterministically encoded by the scaffold cannot count as discovered
from interaction history merely because a learned network reproduces it.

\subtopic{47.2 Formation and certification split}

The same sample must not be used both to select a workspace and to certify its
coherence, causal efficacy, or relational irreducibility. The required split is
\[
\mathcal D
=
\mathcal D_{\mathrm{form}}
\sqcup
\mathcal D_{\mathrm{cal}}
\sqcup
\mathcal D_{\mathrm{cert}}.
\]
The formation sample selects \(\widehat\theta\). The calibration sample fixes
thresholds, regularization strengths, and intervention magnitudes. The
certification sample evaluates the final gates without further structural
selection. In a multi-environment study, the stronger requirement is
\[
\mathcal E_{\mathrm{form}}
\cap
\mathcal E_{\mathrm{cert}}
=
\varnothing,
\]
so that the learned organization is tested under environments not used to
construct it.

\begin{definition}[Formation leakage]\label{def:formation-leakage}
A certification variable \(Y_{\mathrm{cert}}\) leaks into formation when the
conditional law of \(\widehat\theta\) changes under an intervention on
\(Y_{\mathrm{cert}}\) that should be inaccessible during formation:
\[
\mathcal L
\left(
\widehat\theta
\mid
\operatorname{do}(Y_{\mathrm{cert}}=y)
\right)
\neq
\mathcal L
\left(
\widehat\theta
\mid
\operatorname{do}(Y_{\mathrm{cert}}=y')
\right).
\]
A valid discovery claim requires a documented information barrier or an audit
showing that this effect is absent up to tolerance.
\end{definition}

\subtopic{47.3 Origin and reality are orthogonal}

The formation provenance of a workspace and its operational certification are
separate coordinates. A fully scaffolded structure may satisfy every
structural and causal workspace gate. A highly endogenous structure may fail
to become coherent or causally useful. Symbolically,
\[
\operatorname{Endogenous}(W)
\centernot\Longrightarrow
\operatorname{SRW}_{\Theta}(W),
\]
\[
\operatorname{SRW}_{\Theta}(W)
\centernot\Longrightarrow
\operatorname{Endogenous}(W).
\]
The theory therefore treats endogeneity as a provenance axis rather than a
higher level on the workspace hierarchy.

\topic{48}{Operational Endogeneity Signature}

The ideal description-complexity expression from the symbolic outline is
retained as a conceptual limit, but raw Kolmogorov complexity is not
computable. A usable theory must declare a model class, a code, and an
intervention family.

\begin{definition}[Code-relative history contribution]\label{def:mdl-endogeneity}
Let \(\mathfrak M_q\) be a declared coding or predictive model class for a
learned component \(q\). Let
\[
L_{\mathfrak M_q}(q\mid Z)
\]
be its held-out codelength given side information \(Z\). Define
\[
\boxed{
 e_q^{\mathrm{MDL}}
 =
 \left[
 \frac{
 L_{\mathfrak M_q}(q\mid\Sigma)
 -
 L_{\mathfrak M_q}(q\mid\Sigma,\mathcal D_{\mathrm{form}})
 }{
 \max\{L_{\mathfrak M_q}(q\mid\Sigma),1\}
 }
 \right]_{0}^{1},
}
\]
where \([x]_0^1=\min\{1,\max\{0,x\}\}\).
\end{definition}

This quantity is explicitly code-relative. It must not be presented as an
observer-independent metaphysical measure. Its role is to test whether the
formation history provides compressive information about the learned
organization beyond the declared scaffold.

\begin{definition}[Causal formation sensitivity]\label{def:causal-formation-sensitivity}
Let \(\iota_h\) and \(\iota_{h'}\) be matched feasible interventions that alter
formation history while preserving the scaffold, randomization law, sample
size, and intervention cost. For a divergence \(D\), define
\[
\boxed{
 e_q^{\mathrm{causal}}
 =
 \mathbb E_{(h,h')\sim\nu_q}
 D\!\left(
 \mathcal L^{\iota_h}(q),
 \mathcal L^{\iota_{h'}}(q)
 \right).
}
\]
A common null uses temporally or relationally shuffled histories that preserve
local marginals but destroy the candidate formation signal.
\end{definition}

\begin{proposition}[Reparameterization invariance of causal provenance]
\label{prop:causal-provenance-invariance}
Let \(\phi\) be a measurable bijection on the range of \(q\), and let the
divergence \(D\) be invariant under measurable bijections. Then
\[
 e_{\phi(q)}^{\mathrm{causal}}
 =
 e_q^{\mathrm{causal}}.
\]
\end{proposition}

\begin{proof}
Every intervention-induced law of \(\phi(q)\) is the pushforward of the
corresponding law of \(q\). Bijection invariance of \(D\) preserves each
pairwise divergence, and therefore their design-weighted expectation.
\end{proof}

The causal coordinate is therefore insensitive to a mere renaming of the
learned structure. The MDL coordinate need not be invariant unless the coding
class is transformed with the reparameterization.

\subtopic{48.1 Scaffold dependence and cross-scaffold invariance}

Endogeneity does not require independence from all scaffolds. Every learning
system has an architecture, objective, and resource budget. The relevant
question is whether the learned organization depends on the particular
solution-bearing details of the scaffold.

For a family \(\mathfrak S\) of admissible scaffold variants, define the
cross-scaffold identity defect
\[
\boxed{
\delta_q^{\mathrm{scaf}}
=
\inf_{\{J_{\sigma\sigma'}\}}
\mathbb E_{\sigma,\sigma'\sim\pi}
 d_q
\left(
J_{\sigma\sigma'}q_{\sigma},q_{\sigma'}
\right),
}
\]
where \(J_{\sigma\sigma'}\) transports the learned object between admissible
coordinate systems. A low value indicates that the same abstract
organization is recovered under multiple non-solution-bearing scaffolds.

The complete formation signature is
\[
\boxed{
\mathbf E_{\mathrm{form}}(W)
=
\left(
 e_V,e_E,e_H,e_P,e_{\mathcal O},e_B;
 \delta^{\mathrm{scaf}},
 \delta^{\mathrm{leak}}
\right),
}
\]
where every \(e_q\) is reported as the pair
\[
 e_q
=
\left(
 e_q^{\mathrm{MDL}},
 e_q^{\mathrm{causal}}
\right).
\]
No scalar average is permitted to hide a zero history contribution for one of
the structures claimed to have been discovered.

\subtopic{48.2 Formation certificate}

For thresholds \(\Theta_{\mathrm{form}}\), define
\[
\operatorname{Form}_{\Theta_{\mathrm{form}}}(W)=1
\]
only when all applicable gates pass:
\[
\begin{aligned}
& e_q^{\mathrm{MDL}}\geq \tau_q^{\mathrm{MDL}},
&& e_q^{\mathrm{causal}}\geq \tau_q^{\mathrm{causal}},\\
& \delta_q^{\mathrm{scaf}}\leq\varepsilon_q^{\mathrm{scaf}},
&& \delta^{\mathrm{leak}}\leq\varepsilon_{\mathrm{leak}},\\
& \operatorname{SRW}_{\Theta}(W)=1,
&& \operatorname{Pass}(N_{\mathrm{form}})=1.
\end{aligned}
\]
The final workspace gate is included because endogenous discovery of an
unstable or causally irrelevant pattern is not endogenous workspace formation.

\topic{49}{Learned Transfer Structure and Formation Classes}

\subtopic{49.1 Gated membership and edge learning}

Let
\[
z_i\in\{0,1\},
\qquad
a_e\in[0,1]
\]
be node-membership and edge-activation variables satisfying
\[
a_e\leq z_{s(e)},
\qquad
a_e\leq z_{t(e)}.
\]
A weighted learned discrepancy operator is
\[
(D_{a,H}c)_e
=
\sqrt{w_ea_e}
\left(
 c_{t(e)}-H_ec_{s(e)}
\right).
\]
One generic formation objective is
\[
\boxed{
\begin{aligned}
\widehat\theta
\in
\argmin_{z,a,H,P,\mathcal O,B}
\Bigl
\{&
\widehat{\mathcal R}_{\mathrm{task}}
+
\lambda_{\mathrm{coh}}\|D_{a,H}c\|^2
+
\lambda_E\|a\|_1
+
\lambda_V\|z\|_1\\
&+
\lambda_{\mathrm{stab}}\mathcal J_{\mathrm{stab}}
+
\lambda_{\mathrm{leak}}\mathcal J_{\mathrm{leak}}
\Bigr\}.
\end{aligned}
}
\]
The objective is schematic rather than universal. Its purpose is to expose
which quantities were fixed and which were optimized. If \(z\) and \(a\) are
fixed by \(\Sigma\), only the transfer maps or local channels can be claimed as
learned.

\subtopic{49.2 Gauge-identifiable transfer learning}

A learned transfer map is not absolutely identifiable when local coordinates
are latent. Under local invertible changes
\[
X_u'=g_uX_u,
\qquad
X_v'=g_vX_v,
\]
the same abstract transfer is represented by
\[
H_{uv}'=g_vH_{uv}g_u^{-1}.
\]
The identifiable object is therefore the gauge class
\[
[H]_{\mathcal G}
=
\left\{
(g_{t(e)}H_eg_{s(e)}^{-1})_{e\in E}
:
(g_i)_{i\in V}\in\mathcal G
\right\}.
\]
Any formation claim that compares raw matrices across independently trained
systems without gauge alignment is ill posed.

\subtopic{49.3 Transfer identification theorem \(T_{15}\) }

\begin{theorem}[Population transfer identification]\label{thm:transfer-identification}
Let \(X\in\mathbb R^{d_u}\) and \(Y\in\mathbb R^{d_v}\) satisfy
\[
Y=H_*X+\xi,
\qquad
\mathbb E[\xi X^{\mathsf T}]=0,
\]
and assume
\[
\Sigma_{XX}
=
\mathbb E[XX^{\mathsf T}]
\succeq
\kappa I
\]
for some \(\kappa>0\). Then \(H_*\) is the unique minimizer of
\[
\mathcal R(H)
=
\mathbb E\|Y-HX\|_2^2,
\]
with
\[
\boxed{
H_*
=
\Sigma_{YX}\Sigma_{XX}^{-1}.
}
\]
Moreover,
\[
\boxed{
\mathcal R(H)-\mathcal R(H_*)
=
\operatorname{tr}
\left(
(H-H_*)\Sigma_{XX}(H-H_*)^{\mathsf T}
\right)
\geq
\kappa\|H-H_*\|_F^2.
}
\]
\end{theorem}

\begin{proof}
Substitute \(Y=H_*X+\xi\) into the risk. The cross term vanishes because
\(\mathbb E[\xi X^{\mathsf T}]=0\), giving
\[
\mathcal R(H)
=
\mathcal R(H_*)
+
\mathbb E\|(H-H_*)X\|_2^2.
\]
The remaining expectation is the displayed trace. Positive definiteness of
\(\Sigma_{XX}\) gives uniqueness and the lower bound. The normal equations
give the covariance formula.
\end{proof}

\begin{corollary}[Deterministic empirical recovery bound]
\label{cor:empirical-transfer-bound}
For samples \((X_r,Y_r)_{r=1}^n\), write
\[
\widehat\Sigma_{XX}
=
\frac1n\sum_{r=1}^nX_rX_r^{\mathsf T},
\qquad
\widehat M
=
\frac1n\sum_{r=1}^n\xi_rX_r^{\mathsf T}.
\]
If \(\widehat\Sigma_{XX}\succeq\widehat\kappa I\), then ordinary least squares
satisfies
\[
\boxed{
\|\widehat H-H_*\|_F
\leq
\frac{\|\widehat M\|_F}{\widehat\kappa}.
}
\]
For ridge parameter \(\tau>0\),
\[
\boxed{
\|\widehat H_\tau-H_*\|_F
\leq
\frac{
\|\widehat M\|_F+\tau\|H_*\|_F
}{
\widehat\kappa+\tau
}.
}
\]
\end{corollary}

\begin{proof}
The empirical normal equations give
\[
\widehat H-H_*
=
\widehat M\widehat\Sigma_{XX}^{-1}.
\]
For ridge regression,
\[
\widehat H_\tau-H_*
=
(\widehat M-\tau H_*)
(\widehat\Sigma_{XX}+\tau I)^{-1}.
\]
Taking Frobenius and operator norms yields both estimates.
\end{proof}

\begin{proposition}[Gauge covariance of least squares]
\label{prop:ols-gauge-covariance}
Under invertible coordinate transformations
\[
X'=g_uX,
\qquad
Y'=g_vY,
\]
the population least-squares solution transforms as
\[
H_*'=g_vH_*g_u^{-1}.
\]
\end{proposition}

\begin{proof}
The transformed covariances are
\[
\Sigma_{Y'X'}=g_v\Sigma_{YX}g_u^{\mathsf T},
\qquad
\Sigma_{X'X'}=g_u\Sigma_{XX}g_u^{\mathsf T}.
\]
Substitution into the covariance formula and inversion of the second factor
give the result.
\end{proof}

The theorem establishes identification only under excitation and exogeneity.
When \(\Sigma_{XX}\) is singular, the transfer is identified only on the
excited source subspace. When \(\mathbb E[\xi X^{\mathsf T}]\neq0\), regression
recovers a predictive map rather than necessarily the intended causal
transfer.

\subtopic{49.4 Held-out structure-selection theorem \(T_{16}\) }

The same finite-margin argument applies to node membership, edge selection,
boundary selection, or a finite family of gauge-normalized transfer models.

\begin{theorem}[Uniform structure-selection recovery]
\label{thm:uniform-structure-selection}
Let \(\mathfrak T\) be a finite candidate family with population score
\(J:\mathfrak T\to\mathbb R\). Let
\[
\theta_*
=
\argmin_{\theta\in\mathfrak T}J(\theta)
\]
be unique, with margin
\[
\Delta_J
=
\min_{\theta\neq\theta_*}
\bigl(J(\theta)-J(\theta_*)\bigr)>0.
\]
If a held-out estimate \(\widehat J\) satisfies
\[
\sup_{\theta\in\mathfrak T}
|\widehat J(\theta)-J(\theta)|
\leq\varepsilon,
\]
then every empirical minimizer \(\widehat\theta\) obeys
\[
\boxed{
J(\widehat\theta)-J(\theta_*)
\leq2\varepsilon.
}
\]
If \(2\varepsilon<\Delta_J\), then
\[
\boxed{\widehat\theta=\theta_*.}
\]
\end{theorem}

\begin{proof}
By empirical optimality and uniform approximation,
\[
J(\widehat\theta)
\leq
\widehat J(\widehat\theta)+\varepsilon
\leq
\widehat J(\theta_*)+\varepsilon
\leq
J(\theta_*)+2\varepsilon.
\]
If \(\widehat\theta\neq\theta_*\), the margin would force an excess larger than
\(2\varepsilon\), a contradiction.
\end{proof}

The theorem does not guarantee that the candidate class contains the correct
organization. It states exactly what follows once a finite class, score, and
uniform certification error are declared.

\subtopic{49.5 Formation classes}

The provenance classes are now defined relative to
\[
(\Sigma,\mathcal L,\mathcal I_{\mathrm{form}},\mathfrak M).
\]
They are not absolute labels.

\begin{description}[leftmargin=2.1em,style=nextline]
\item[Exogenously scaffolded.]
The claimed nodes, edges, transfer maps, channels, operations, or boundary are
recoverable from the scaffold alone at the declared tolerance. Their history
advantages and causal formation sensitivities vanish.

\item[Adaptively parameterized.]
The membership and channel topology are fixed, while transfer parameters,
local operators, or readouts are learned from interaction history. Typically
\[
e_H+e_P+e_{\mathcal O}>0,
\qquad
e_V\approx e_E\approx e_B\approx0.
\]

\item[Structurally adaptive.]
At least one of membership, edges, or boundary is selected from history, but
the search family and objective remain externally specified.

\item[Endogenously constituted relative to the scaffold class.]
Membership, transfer organization, and operational boundary have nonzero
history contribution and causal formation sensitivity, survive cross-scaffold
transport, pass leakage and matched-null tests, and produce a certified
workspace on held-out environments.
\end{description}

The phrase ``relative to the scaffold class'' is mandatory. No finite
experiment can establish independence from every possible external influence.
